%==============================================================================
%  Appendix
%  Each appendix becomes a "chapter" under the \appendix block in main.tex
%  (so they appear as Appendix A, B, C, ...).
%
%  This thesis uses a documentary design (no interviews), so the appendices
%  document the corpus, the codebooks and supplementary analytical material.
%==============================================================================

\chapter{Document corpus}
\label{app:doc-corpus}

\section*{Analysed corpus}

The corpus consists of a single document, listed below. Every coded
passage in \texttt{Readthrough\_Notes\_v1.md} is drawn from this
document.

\begin{longtable}{@{} l p{6.5cm} l l l @{}}
\toprule
\textbf{ID} & \textbf{Title} & \textbf{Issuer} & \textbf{Year} & \textbf{Type} \\
\midrule
\endfirsthead
\toprule
\textbf{ID} & \textbf{Title} & \textbf{Issuer} & \textbf{Year} & \textbf{Type} \\
\midrule
\endhead
\bottomrule
\endfoot
D01 & EIC Work Programme 2026 (Pathfinder \& STEP Scale Up sections, C(2025) 7410) & European Commission & 2026 & Policy doc \\
\end{longtable}

\section*{Legal instruments referenced within the corpus}

The two regulations below are cited by the Work Programme as the
binding legal basis for the rules it implements. They are listed here
for reference only --- they were not themselves coded or analysed, and
they appear in the thesis only where the Work Programme cites them.

\begin{longtable}{@{} l p{6.5cm} l l l @{}}
\toprule
\textbf{ID} & \textbf{Title} & \textbf{Issuer} & \textbf{Year} & \textbf{Type} \\
\midrule
\endfirsthead
\toprule
\textbf{ID} & \textbf{Title} & \textbf{Issuer} & \textbf{Year} & \textbf{Type} \\
\midrule
\endhead
\bottomrule
\endfoot
R01 & STEP Regulation (EU) 2024/795 & European Parliament \& Council & 2024 & Regulation \\
R02 & Establishing Horizon Europe --- Regulation (EU) 2021/695 & European Parliament \& Council & 2021 & Regulation \\
\end{longtable}

%==============================================================================
\chapter{Deductive codebook}
\label{app:codebook-deductive}

This appendix reproduces the deductive codebook applied in Pass~1 of the
analysis. The codebook is organised in six sections that follow the
theoretical framework set out in Chapter~\ref{ch:background}: the three
pillars of \textcite{scott2014institutions}, the four institutional
logics of \textcite{thornton2012institutional} as developed from
\textcite{friedland1991bringing}, the three types of isomorphism in
\textcite{dimaggio1983iron}, and the legitimacy typology of
\textcite{suchman1995managing}. The full v1 codebook is held in the
project repository at \texttt{04\_codebook/output/Codebook\_v1\_Deductive.md}.

\paragraph{Unit of analysis.} The sentence or clause, extended to the
full paragraph only where the meaning cannot be extracted from a
shorter unit. A passage may receive multiple codes from different
sections simultaneously; codes from each section are not mutually
exclusive.

\paragraph{Programme tag.} Each coded passage is marked \texttt{[PATH]}
(Pathfinder), \texttt{[STEP]} (STEP Scale Up), or \texttt{[SHARED]}
for introductory or shared sections.

\paragraph{Ambiguity flag.} Uncertain coding decisions are marked
\texttt{[?]} during Pass~1 and resolved in Pass~2.

\section*{Section 1 --- Regulative pillar}

Rules, laws, monitoring, and sanctions. Legitimacy through compliance.
Ask: what must actors do, avoid, or demonstrate in order to be admitted
and remain in good standing?

\begin{longtable}{@{} l p{4.3cm} p{6.3cm} @{}}
\toprule
\textbf{Code} & \textbf{Definition} & \textbf{Anchor / example} \\
\midrule
\endfirsthead
\toprule
\textbf{Code} & \textbf{Definition} & \textbf{Anchor / example} \\
\midrule
\endhead
\bottomrule
\endfoot
\texttt{REG\_TRL} & TRL as formal eligibility threshold or scope boundary. & \enquote{Technology Readiness Levels from 1 to 4}; activities beyond TRL~4 not eligible. \\
\texttt{REG\_ELIGIBILITY} & Formal admissibility: who may apply and under what conditions, independent of quality. & Concurrent-submission prohibition; nationality and establishment rules. \\
\texttt{REG\_SCOPE} & State-directed thematic channelling toward pre-defined technology themes or societal challenges. & Named Pathfinder Challenge topics; STEP priority domains. \\
\texttt{REG\_IP} & Intellectual-property conditions on funded entities (ownership, transfer, EU retention). & EIC Fund ensuring companies keep IP in the EU; Annex IP provisions. \\
\texttt{REG\_SANCTION} & Formal consequences for non-compliance, failure or rule violation (exclusion, clawback, termination). & Three-strike exclusion after three unsuccessful STEP submissions. \\
\texttt{REG\_MANDATE} & Programme provisions explicitly grounded in named EU legislative instruments. & STEP Regulation (EU)~2024/795; Decision C(2025)\,7410. \\
\texttt{REG\_FINREG} & Financial-sector regulation imported into innovation funding. & KYC / AML compliance by the EIC Fund or EIF. \\
\end{longtable}

\section*{Section 2 --- Normative pillar}

Values, norms, role expectations, and obligations. Legitimacy through
appropriateness. Ask: what kind of actor should an applicant be, and
what obligations does participation entail?

\begin{longtable}{@{} l p{4.3cm} p{6.3cm} @{}}
\toprule
\textbf{Code} & \textbf{Definition} & \textbf{Anchor / example} \\
\midrule
\endfirsthead
\toprule
\textbf{Code} & \textbf{Definition} & \textbf{Anchor / example} \\
\midrule
\endhead
\bottomrule
\endfoot
\texttt{NORM\_IDENTITY} & Construction of the ideal applicant identity through aspirational or prescriptive language. & \enquote{Do you have an ambitious vision \dots\,?}; \enquote{game-changing innovative technology}. \\
\texttt{NORM\_OBLIGATION\_SOCIETAL} & Research framed as carrying a duty to society beyond commercial success. & \enquote{Respond to societal needs and / or provide solutions for global challenges}. \\
\texttt{NORM\_OBLIGATION\_STRATEGIC} & Participation framed as serving European strategic interests and competitiveness. & \enquote{Propelling Europe's economic, industrial, and technological competitiveness}. \\
\texttt{NORM\_COLLABORATION} & Interdisciplinarity and cross-sector collaboration valorised as a signal of legitimacy. & \enquote{Interdisciplinary team of researchers and innovators}. \\
\texttt{NORM\_DEI} & Diversity / gender requirements with evaluative weight, not boilerplate. & \enquote{Adequate gender balance, with a credible plan to fill the gaps}. \\
\texttt{NORM\_SMARTMONEY} & EIC relationship framed as VC-style value-add partnership beyond the grant. & \enquote{Smart money}; Business Acceleration Services. \\
\end{longtable}

\section*{Section 3 --- Cognitive pillar}

Shared schemas, taken-for-granted assumptions, and naturalised
categories. Legitimacy through conformity with what is self-evident.
Ask: what does the programme treat as obvious, unquestioned or natural
without arguing for it?

\begin{longtable}{@{} l p{4.3cm} p{6.3cm} @{}}
\toprule
\textbf{Code} & \textbf{Definition} & \textbf{Anchor / example} \\
\midrule
\endfirsthead
\toprule
\textbf{Code} & \textbf{Definition} & \textbf{Anchor / example} \\
\midrule
\endhead
\bottomrule
\endfoot
\texttt{COG\_DEEPTECH} & \enquote{Deep tech} used as a self-evident category without definition. & \enquote{Earliest stages of scientific, technological or deep tech research}. \\
\texttt{COG\_LINEAR} & Linear innovation model (basic $\rightarrow$ applied $\rightarrow$ market) treated as a self-evident schema. & TRL~1--4 $\rightarrow$ 5--8 $\rightarrow$ commercial as natural progression. \\
\texttt{COG\_DISRUPTION} & Disruption treated as the self-evidently appropriate goal of deep-tech research. & \enquote{Disrupt a field and a market or create new opportunities}. \\
\texttt{COG\_THREAT} & European strategic vulnerability treated as established fact, not argued. & \enquote{Strategic dependencies of the Union} as background premise. \\
\texttt{COG\_MARKETFAILURE} & Private market failure treated as objective premise requiring no demonstration. & \enquote{Market gap in financing deep tech scale up companies in Europe}. \\
\texttt{COG\_FINANCIALISATION} & Equity investment treated as the self-evidently appropriate scale-up mechanism. & Equity-only STEP instrument; funding-round catalysis as a natural goal. \\
\end{longtable}

\section*{Section 4 --- Institutional logics}

Cross-cutting codes identifying which deep cultural ordering principle
organises a passage's assumptions about what matters, who has
authority, and what counts as success.

\begin{longtable}{@{} l p{4.3cm} p{6.3cm} @{}}
\toprule
\textbf{Code} & \textbf{Definition} & \textbf{Anchor / example} \\
\midrule
\endfirsthead
\toprule
\textbf{Code} & \textbf{Definition} & \textbf{Anchor / example} \\
\midrule
\endhead
\bottomrule
\endfoot
\texttt{LOGIC\_SCIENCE} & Excellence, peer review, discovery-driven research, knowledge production as an end in itself. & Pathfinder Open evaluated on scientific merit; \enquote{validate the scientific basis}. \\
\texttt{LOGIC\_MARKET} & Commercial viability, investor returns, competition, demand signals, growth as the measure of value. & Revenue and market opportunity as criteria; \enquote{first come, first served} allocation. \\
\texttt{LOGIC\_STATE} & Sovereignty, industrial strategy, national security, the state as ultimate arbiter of strategic value. & STEP Regulation as mandate; Sovereignty Seal; IP retention in the EU. \\
\texttt{LOGIC\_PROFESSION} & Credentialled expertise, professional standards, the authority of specialised practitioners. & Expert jury; Technology Due Diligence; FTO analysis. \\
\texttt{LOGIC\_TENSION} & Two or more logics simultaneously present in observable conflict, surface-visible in the text. & Pathfinder Open (science) and Challenges (state) as co-existing sub-programmes. \\
\end{longtable}

\texttt{LOGIC\_TENSION} is always applied alongside the substantive
logic codes it links. It is reserved for tension legible in the text
itself; tensions constructed across passages by analytical argument are
recorded as Pattern~A observations (see
Appendix~\ref{app:codebook-inductive}).

\section*{Section 5 --- Isomorphic mechanisms}

Codes identifying how the programme produces conformity and from what
source that pressure originates.

\begin{longtable}{@{} l p{4.3cm} p{6.3cm} @{}}
\toprule
\textbf{Code} & \textbf{Definition} & \textbf{Anchor / example} \\
\midrule
\endfirsthead
\toprule
\textbf{Code} & \textbf{Definition} & \textbf{Anchor / example} \\
\midrule
\endhead
\bottomrule
\endfoot
\texttt{ISO\_COERCIVE} & Conformity through legal mandate, regulatory requirement, or the formal authority of a more powerful actor. & STEP Regulation; IP retention clause; three-strike rule; KYC / AML. \\
\texttt{ISO\_MIMETIC} & Conformity through imitation of a perceived successful external model. & Strategic Goal 4 benchmarking against the USA and Asia; adoption of structural features characterised in the developmental-state literature as comparator-template mechanics; VC jury and pitch-deck vocabulary. \\
\texttt{ISO\_NORMATIVE} & Conformity through professional norms, standards, and best practices. & Business plan and pitch deck requirements; mandatory FTO analysis; DEI norms. \\
\end{longtable}

\section*{Section 6 --- Legitimacy strategies}

Codes identifying how the programme constructs and claims legitimacy,
both for itself and for the innovation activities it funds, following
\textcite{suchman1995managing}.

\begin{longtable}{@{} l p{4.3cm} p{6.3cm} @{}}
\toprule
\textbf{Code} & \textbf{Definition} & \textbf{Anchor / example} \\
\midrule
\endfirsthead
\toprule
\textbf{Code} & \textbf{Definition} & \textbf{Anchor / example} \\
\midrule
\endhead
\bottomrule
\endfoot
\texttt{LEG\_PRAGMATIC} & Justification through concrete utility to identifiable stakeholders. & EUR~10--30\,M investment; EUR~50--150\,M funding-round target; BAS provisions. \\
\texttt{LEG\_MORAL} & Justification through alignment with broader social values or shared normative commitments. & Sovereignty and strategic autonomy framed as obligation; societal-challenge framing. \\
\texttt{LEG\_COGNITIVE} & Existence or design presented as natural, inevitable, or self-evidently necessary. & Programme naturalised by the European tech gap; deep tech as a natural category. \\
\texttt{LEG\_SEAL} & Sovereignty (STEP) Seal as a portable, transferable legitimacy artefact. & Seal awarded to NO~GO proposals meeting strategic criteria; valid for use with InvestEU and national programmes. \\
\end{longtable}

\section*{Code index}

Thirty-one codes in total across the six sections. The right-hand
column records the programme(s) in which the code chiefly applies.

\begin{longtable}{@{} l l l @{}}
\toprule
\textbf{Code} & \textbf{Section} & \textbf{Primary programme} \\
\midrule
\endfirsthead
\toprule
\textbf{Code} & \textbf{Section} & \textbf{Primary programme} \\
\midrule
\endhead
\bottomrule
\endfoot
\texttt{REG\_TRL} & Regulative & PATH \\
\texttt{REG\_ELIGIBILITY} & Regulative & Both \\
\texttt{REG\_SCOPE} & Regulative & Both \\
\texttt{REG\_IP} & Regulative & Both (stronger in STEP) \\
\texttt{REG\_SANCTION} & Regulative & Both (stronger in STEP) \\
\texttt{REG\_MANDATE} & Regulative & STEP \\
\texttt{REG\_FINREG} & Regulative & STEP \\
\texttt{NORM\_IDENTITY} & Normative & Both \\
\texttt{NORM\_OBLIGATION\_SOCIETAL} & Normative & PATH \\
\texttt{NORM\_OBLIGATION\_STRATEGIC} & Normative & STEP \\
\texttt{NORM\_COLLABORATION} & Normative & PATH \\
\texttt{NORM\_DEI} & Normative & Both \\
\texttt{NORM\_SMARTMONEY} & Normative & STEP \\
\texttt{COG\_DEEPTECH} & Cognitive & Both \\
\texttt{COG\_LINEAR} & Cognitive & PATH \\
\texttt{COG\_DISRUPTION} & Cognitive & Both \\
\texttt{COG\_THREAT} & Cognitive & STEP \\
\texttt{COG\_MARKETFAILURE} & Cognitive & STEP \\
\texttt{COG\_FINANCIALISATION} & Cognitive & STEP \\
\texttt{LOGIC\_SCIENCE} & Logics & PATH \\
\texttt{LOGIC\_MARKET} & Logics & Both \\
\texttt{LOGIC\_STATE} & Logics & STEP \\
\texttt{LOGIC\_PROFESSION} & Logics & Both \\
\texttt{LOGIC\_TENSION} & Logics & Both \\
\texttt{ISO\_COERCIVE} & Isomorphism & STEP \\
\texttt{ISO\_MIMETIC} & Isomorphism & Both \\
\texttt{ISO\_NORMATIVE} & Isomorphism & Both \\
\texttt{LEG\_PRAGMATIC} & Legitimacy & Both \\
\texttt{LEG\_MORAL} & Legitimacy & Both \\
\texttt{LEG\_COGNITIVE} & Legitimacy & Both \\
\texttt{LEG\_SEAL} & Legitimacy & STEP \\
\end{longtable}

%==============================================================================
\chapter{Institutional logics quick-reference matrix}
\label{app:logics-matrix}

This appendix reproduces the quick-reference matrix used during coding
to discipline the application of the four institutional-logic codes
(\texttt{LOGIC\_SCIENCE}, \texttt{LOGIC\_MARKET}, \texttt{LOGIC\_STATE},
\texttt{LOGIC\_PROFESSION}). It compares the logics across the
attributes that mattered in practice for coding decisions: the central
question each logic asks, its source of authority, its measure of
success, who decides, the typical vocabulary in which it surfaces, the
instrument it favours, the kind of applicant it expects, and the
signal it gives off in the Pathfinder and STEP texts respectively. The
matrix is the operational counterpart to the conceptual treatment in
\cref{sec:bg-it-logics} and supplements the logic codes formalised in
\cref{app:codebook-deductive}, Section~4. The basis is
\textcite{friedland1991bringing} and \textcite{thornton2012institutional}.

\begin{longtable}{@{} p{2.4cm} p{2.6cm} p{2.6cm} p{2.6cm} p{2.6cm} @{}}
\toprule
& \textbf{\texttt{LOGIC\_SCIENCE}} & \textbf{\texttt{LOGIC\_MARKET}} & \textbf{\texttt{LOGIC\_STATE}} & \textbf{\texttt{LOGIC\_PROFESSION}} \\
\midrule
\endfirsthead
\toprule
& \textbf{\texttt{LOGIC\_SCIENCE}} & \textbf{\texttt{LOGIC\_MARKET}} & \textbf{\texttt{LOGIC\_STATE}} & \textbf{\texttt{LOGIC\_PROFESSION}} \\
\midrule
\endhead
\bottomrule
\endfoot
{\footnotesize\textbf{Central question}} & {\footnotesize Is this excellent science?} & {\footnotesize Will this create and capture market value?} & {\footnotesize Does this serve European strategic interests?} & {\footnotesize Does this meet professional standards?} \\
\addlinespace
{\footnotesize\textbf{Source of authority}} & {\footnotesize Scientific expertise and peer judgement} & {\footnotesize Investors and market signals} & {\footnotesize Legislative mandate and public institutions} & {\footnotesize Credentialled practitioners and professional bodies} \\
\addlinespace
{\footnotesize\textbf{Measure of success}} & {\footnotesize Novelty, rigour, frontier contribution} & {\footnotesize Commercial viability, growth, return} & {\footnotesize Sovereignty, security, geopolitical positioning} & {\footnotesize Conformity with professional norms and expert endorsement} \\
\addlinespace
{\footnotesize\textbf{Who decides}} & {\footnotesize Peer reviewers and scientific jury} & {\footnotesize Co-investors and market demand} & {\footnotesize EU institutions and the STEP Regulation} & {\footnotesize External specialists (IP, technology, investment)} \\
\addlinespace
{\footnotesize\textbf{Typical vocabulary}} & {\footnotesize \enquote{scientific basis}, \enquote{proof of principle}, \enquote{frontier}, \enquote{novel}, \enquote{interdisciplinary}, \enquote{TRL~1--4}} & {\footnotesize \enquote{market opportunity}, \enquote{commercial potential}, \enquote{funding round}, \enquote{investor}, \enquote{scale}, \enquote{revenue}, \enquote{first come, first served}} & {\footnotesize \enquote{strategic dependencies}, \enquote{sovereignty}, \enquote{European interests}, \enquote{protect}, \enquote{strategic autonomy}, \enquote{geopolitical}} & {\footnotesize \enquote{due diligence}, \enquote{expert}, \enquote{leading specialists}, \enquote{Freedom to Operate}, \enquote{credible plan}, \enquote{assessment}} \\
\addlinespace
{\footnotesize\textbf{Favoured instrument}} & {\footnotesize Open grant, peer-reviewed call} & {\footnotesize Equity, co-investment, competitive pitch} & {\footnotesize Regulated mandate, conditioned award, IP clause} & {\footnotesize Expert-validated selection, professional certification} \\
\addlinespace
{\footnotesize\textbf{Expected applicant}} & {\footnotesize Rigorous researcher with an ambitious scientific vision} & {\footnotesize Commercially viable company with a credible growth plan} & {\footnotesize Strategic asset with European anchoring} & {\footnotesize Professionally organised entity able to speak expert languages} \\
\addlinespace
{\footnotesize\textbf{Pathfinder signal}} & {\footnotesize Primary in Pathfinder Open} & {\footnotesize Secondary --- commercialisation pathway required} & {\footnotesize Tertiary --- visible in Pathfinder Challenges} & {\footnotesize Present in jury evaluation and expert review} \\
\addlinespace
{\footnotesize\textbf{STEP signal}} & {\footnotesize Residual --- \enquote{deep tech in nature} excellence criterion} & {\footnotesize Primary in instrument design (equity, rounds, co-investment)} & {\footnotesize Primary in programme purpose (sovereignty, strategic tech)} & {\footnotesize Primary in due diligence and evaluation architecture} \\
\end{longtable}

\section*{Logic-tension pairs}

The pairs below name the tensions that recur across the corpus when
two logics are claimed in a single passage. The pair vocabulary is
used in Pattern~A (\cref{sec:app-c-patterns}) and informs the
surface-visibility standard for \texttt{LOGIC\_TENSION} (Rule~1,
\cref{sec:app-c-rules}).

\begin{longtable}{@{} p{3.0cm} p{10.0cm} @{}}
\toprule
\textbf{Tension pair} & \textbf{What it looks like in the text} \\
\midrule
\endfirsthead
\toprule
\textbf{Tension pair} & \textbf{What it looks like in the text} \\
\midrule
\endhead
\bottomrule
\endfoot
SCIENCE~$\leftrightarrow$~STATE & Scientific-merit criterion applied within a state-defined thematic scope; researcher autonomy constrained by strategic priorities. \\
MARKET~$\leftrightarrow$~STATE & Market instrument (equity) used to achieve state goal (sovereignty, IP retention); commercial logic form, geopolitical logic content. \\
SCIENCE~$\leftrightarrow$~MARKET & Commercialisation pathway required from TRL~1--4 research; impact criteria imposed on basic science. \\
STATE~$\leftrightarrow$~PROFESSION & Geopolitical priority set by regulation but validated by expert due diligence; tension between political authority and expert authority. \\
\end{longtable}

%==============================================================================
\chapter{Inductive codes and pattern-level observations}
\label{app:codebook-inductive}

Pass~2 added no new codes to the deductive scheme: the thirty-one v1
codes carry forward unchanged. The inductive update operates instead at
two levels above the individual code. First, seven coding-rule
clarifications discipline how the v1 codes are applied at the margin
(\cref{sec:app-c-rules}). Second, six pattern-level observations
(\cref{sec:app-c-patterns}) name recurring configurations across the
corpus that no single passage can carry on its own. A seventh candidate
template was tested and not confirmed; the negative finding is recorded
in \cref{sec:app-c-not-confirmed}. The full inductive working codebook
is held at \texttt{04\_codebook/output/Codebook\_v2\_Inductive.md}.

\section{Pass~2 coding-rule clarifications}
\label{sec:app-c-rules}

The seven rules below were introduced or codified during Pass~2 to
discipline coding decisions at the margin. Each rule refines how an
existing v1 code applies; no rule changes a code's definition.

\begin{description}[leftmargin=*, style=nextline]
  \item[Rule~1 --- Surface-visibility standard for \texttt{LOGIC\_TENSION}.]
    Tension must be legible in the text itself. Tensions constructed by
    analytical argument across passages are recorded as Pattern~A
    observations, not as \texttt{LOGIC\_TENSION} codings.
  \item[Rule~2 --- STEP catalytic-investment co-coding (Pattern~E).]
    In STEP passages describing the catalytic mechanism, the
    market-interest eligibility criterion, or the market-failure award
    criterion, code \texttt{LOGIC\_STATE}. Add \texttt{LOGIC\_MARKET}
    as a co-code only where market mechanics (equity instrument,
    case-by-case negotiation, co-investment structure) are
    operationally surface-visible in that passage. The structural
    circularity is recorded as Pattern~E.
  \item[Rule~3 --- Contractual discretionary coercion (Pattern~B).]
    \texttt{ISO\_COERCIVE} extends to contractual discretionary
    authority exercised by the EIC under the grant agreement, where
    the grant agreement is itself a legal instrument. Code where the
    authority asserted carries explicit sanction force (termination,
    suspension, mandated amendment) and the source of obligation is
    the grant agreement or a Programme Manager acting under it.
    Co-code \texttt{LOGIC\_STATE} where the state is named as the
    actor wielding that power. Soft-direction language without
    sanction force (\enquote{may be requested to join},
    \enquote{encouraged to collaborate}) remains below threshold and
    is carried as Pattern~B observation only.
  \item[Rule~4 --- State social-equity template (Pattern~C).]
    Where DEI mandates appear as state-mandated social policy
    (recurrent template language, framing as a legislative
    commitment, embedding in the evaluation or tiebreaker
    architecture, institutional self-monitoring of equity compliance),
    code \texttt{LOGIC\_STATE}. Add \texttt{NORM\_DEI} as a co-code
    only where the requirement also operates as a
    normative-organisational standard with evaluation weight against
    a \enquote{credible plan}-type gate.
  \item[Rule~5 --- Business Acceleration Services (BAS).]
    Where BAS is described with professional-services framing
    (\enquote{leading expertise}, \enquote{externally contracted},
    \enquote{bespoke}), code \texttt{LOGIC\_PROFESSION} as the
    organising logic. Add \texttt{LOGIC\_MARKET} as a co-code only
    where market actors (corporates, investors, ecosystem actors) are
    explicitly named as the pool being mobilised under professional
    delivery; here \texttt{LOGIC\_MARKET} describes the actor pool
    rather than independently governing the passage.
  \item[Rule~6 --- Structural state inference.]
    \texttt{LOGIC\_STATE} requires that the state organises the
    passage as the actor whose institution, architecture, or
    strategic act is the subject. Two cases meet the threshold:
    (a)~state vocabulary or an explicit strategic / geopolitical /
    sovereignty rationale governs the passage; (b)~the institutional
    act being performed (eligibility redrawing, scope determination,
    directive intervention) is itself state-institutional and is the
    subject of the passage. Where the passage's subject is something
    else and the state inference rests on external knowledge, code
    None; the observation is preserved as an analytical note.
  \item[Rule~7 --- Vocabulary versus organising logic.]
    A logic is coded only where it \emph{organises} the passage,
    not where it merely appears as vocabulary. Below-threshold
    categories include VC vocabulary on administrative actions
    (\enquote{follow on investments}, \enquote{Booster},
    \enquote{Fast Track}); state-as-administrative-funder language
    without strategic rationale (\enquote{direct financial support});
    management-register language without invoked credentialled
    authority (\enquote{proactive management}); and
    normative-identity language describing applicant ideal types
    (\enquote{visionary researchers}) rather than science-organising
    governance. Apply alongside rules~1 and~6.
\end{description}

\section{Pattern-level observations}
\label{sec:app-c-patterns}

Six patterns were confirmed in Pass~2. Each is described below as a
named configuration with a brief data structure in
\textcite{gioia2013seeking} style: the entries that anchor it, the
mechanism it captures, and where it appears in
Chapter~\ref{ch:results} and Chapter~\ref{ch:discussion}. Where a
pattern produced a coding-rule refinement, the linked rule is named.

\paragraph{Pattern~A --- Tension-Dissolution.}
Logics are presented as naturally complementary; the text dissolves a
tension that is structurally present but rhetorically suppressed. The
analytical inverse of \texttt{LOGIC\_TENSION}: Pattern~A marks
surface-suppressed complementarity, where \texttt{LOGIC\_TENSION} marks
surface-visible conflict. Linked rule: Rule~1.

\begin{itemize}[leftmargin=*]
  \item \textbf{Anchor entries:} [003], [005], [008], [019], [039],
    [044], [056], [057], [102], [109].
  \item \textbf{Signature:} co-equal listing of science / market / state
    sub-criteria without a conflict marker; juxtaposition of
    \enquote{cross-fertilisation and exploitation}; market vocabulary
    folded into science quality criteria.
  \item \textbf{Reading:} Chapter~\ref{ch:results} as a recurrent
    architectural move; Chapter~\ref{ch:discussion} as a legitimacy
    strategy.
\end{itemize}

\paragraph{Pattern~B --- Coercive Governance Cluster.}
A cluster spanning a graded intensity in which the EIC Programme
Manager holds discretionary authority over capital continuity, roadmap
direction, portfolio placement, and inter-consortium collaboration. The
hard end carries termination authority; the soft end is encouragement.
Linked rule: Rule~3.

\begin{itemize}[leftmargin=*]
  \item \textbf{Anchor entries:} [024]--[027], [061], [115].
  \item \textbf{Intensity gradient:} hard (termination, suspension,
    mandated amendment) $\rightarrow$ intermediate (directive roadmap
    revision under threat) $\rightarrow$ soft (\enquote{may be
    requested}, \enquote{encouraged}).
  \item \textbf{Reading:} Chapter~\ref{ch:results} as a contractual
    architecture; Chapter~\ref{ch:discussion} as coercive isomorphism
    extended into post-award discretion.
\end{itemize}

\paragraph{Pattern~C --- State Social Equity Template.}
Near-identical passages embedding social-equity mandates (gender
balance, female-researcher empowerment, underrepresented groups) recur
across instruments and sections with minimal variation, indicating a
standardised insert. Linked rule: Rule~4.

\begin{itemize}[leftmargin=*]
  \item \textbf{Anchor entries:} [001], [017], [047], [063]; with
    [059] (tiebreaker architecture) connecting the template to
    evaluation ranking.
  \item \textbf{Signature:} \enquote{adequate gender balance, with a
    credible plan to fill the gaps} reproduced verbatim across
    instruments; the same gender ranking appearing above SME
    participation and Member-State geographic coverage in tiebreakers.
  \item \textbf{Reading:} Chapter~\ref{ch:results} as a state insert;
    Chapter~\ref{ch:discussion} as the developmental-state dimension
    of the EIC pursuing collective political interests through an
    innovation instrument.
\end{itemize}

\paragraph{Pattern~D --- Challenge Chapter Logic Succession Template.}
Each Challenge chapter reproduces a consistent logic succession across
its sub-sections, with the succession itself functioning as a template
applied across substantively different technology domains. No linked
rule: Pattern~D is a property of how correctly applied codings
configure across a chapter, not of any single passage.

\begin{itemize}[leftmargin=*]
  \item \textbf{Anchor entries:} [095]--[116] (Challenges~II.2.1
    and~II.2.2); [117]--[129] (Challenge~II.2.3).
  \item \textbf{Template:} Scope / background $\rightarrow$
    \texttt{LOGIC\_STATE}; Specific Objectives $\rightarrow$
    \texttt{LOGIC\_SCIENCE}; Portfolio Approach $\rightarrow$
    \texttt{LOGIC\_PROFESSION}; Expected Outcomes $\rightarrow$
    \texttt{LOGIC\_STATE}.
  \item \textbf{Reading:} Chapter~\ref{ch:results} as a chapter-level
    template; Chapter~\ref{ch:discussion} as a meso-level
    institutional reproduction independent of technological content.
\end{itemize}

\paragraph{Pattern~E --- STEP Circular Market Validation.}
The state constructs the market signal it then requires as proof of
market legitimacy. The 20--33\% EIC anchor position structurally
generates the \enquote{market interest} the programme then evaluates
against. Linked rule: Rule~2.

\begin{itemize}[leftmargin=*]
  \item \textbf{Anchor entries:} [132], [140], [141], [145], [176].
  \item \textbf{Signature:} co-presence of \emph{market interest} as
    eligibility criterion and \emph{market failure} as award
    criterion, combined with the anchor position; equity-only
    instrument with case-by-case negotiation.
  \item \textbf{Reading:} Chapter~\ref{ch:results} as a circular
    architecture; Chapter~\ref{ch:discussion} as a
    developmental-state design that hybridises market mechanics with
    state access architecture.
\end{itemize}

\paragraph{Pattern~F --- TRL as Cross-Pillar Legitimation Artefact.}
TRL functions simultaneously as \texttt{REG\_TRL} (regulative gate
determining eligibility, scope boundaries, budget allocation) and as
\texttt{COG\_LINEAR} (the taken-for-granted scientific-developmental
ladder). The two codes are not in tension; they co-code in the same
passages and together constitute a legitimation mechanism. The
cognitive pillar provides scientific-register surface cover; the
regulative pillar performs the gatekeeping work beneath. No linked
rule: v1 already permits co-coding.

\begin{itemize}[leftmargin=*]
  \item \textbf{Anchor entries:} [012], [044], [046]; recurs across
    Pathfinder evaluation passages.
  \item \textbf{Triple load:} cognitive (scientific-measurement
    register), regulative (eligibility gate), and mimetic (TRL framework
    inherited from US defence and NASA).
  \item \textbf{Reading:} Chapter~\ref{ch:results} as the gate-as-ladder
    mechanism; Chapter~\ref{ch:discussion} as cognitive legitimacy in
    \textcite{suchman1995managing}, connecting to mimetic isomorphism.
\end{itemize}

\section{Candidate pattern tested and not confirmed}
\label{sec:app-c-not-confirmed}

\paragraph{Sustainability as a state template.}
The hypothesis was that sustainability framing recurs across Challenge
chapters as a standardised state insert, analogous to Pattern~C. The
universality test fails at Challenge~II.2.2 (Biotechnology for Healthy
Ageing), where no sustainability or environmental language appears.
Sustainability is subject-matter-dependent, not a universal template.

\paragraph{Substantive finding preserved.}
When sustainability does appear, it enters through one of four
mechanisms: (i)~Commission-Communication policy-citation, where the
Challenge cites a Communication that carries the sustainability frame;
(ii)~explicit environmental-consequence framing as descriptive context;
(iii)~sector-definition, where sustainability is constitutive of the
technology field; (iv)~a technical efficiency-proxy, where
sustainability is absorbed into methodology requirements through
constrained-resource specifications. The contrast with Pattern~C
sharpens what the equity template actually templates: the institutional
architecture carries a small set of state priorities universally
(gender mainstreaming) while letting others (sustainability) enter
through subject-matter-appropriate channels.

%==============================================================================
\chapter{Readthrough audit trail (excerpt)}
\label{app:audit-trail}

This appendix reproduces a representative excerpt from
\texttt{Readthrough\_Notes\_v1.md} to evidence the dependability and
confirmability of the coding procedure. The excerpt comprises four
consecutive entries from the EIC Pathfinder Open chapter,
[043]--[046], chosen because they show, side by side, the four
decision types the readthrough produced: an uncoded entry with strong
science-logic undertones, an uncoded entry with explicit researcher
autonomy, a multi-coded entry with a \texttt{[?]} flag for inductive
review, and a fully coded entry with a surface-visible
\texttt{LOGIC\_TENSION}. The standard entry template (in force from
[011] onward) is shown first; the four entries are then reproduced
verbatim from the working notes. The full notes file is held in the
project repository.

\section*{Standard entry template}

\begin{verbatim}
### [NNN] --- Short descriptive title

**Document location:** [Chapter / Section / Page reference]
**Programme:** [PATH] / [STEP] / [SHARED]

**Source text:**
> *"Exact quoted passage"*

**Question:** What was asked.

**Analysis:** Full reasoning.

**Coding:** LOGIC_X + LOGIC_Y + LOGIC_TENSION

**[?] Flag (if applicable):** Unresolved question for inductive update.
\end{verbatim}

\section*{Entry [043] --- EIC Pathfinder Open: high-risk research paragraph}

\textbf{Document location:} II.~EIC Pathfinder --- II.1~EIC Pathfinder
Open / high-risk paragraph. \\
\textbf{Programme:} \texttt{[PATH]}.

\textbf{Source text:}
\begin{quote}
\enquote{EIC Pathfinder Open may support your work, especially if it is
highly risky: you may set out to try things that will not work; you
may be faced with questions that nobody knows the answer to yet; you
may realise that there are many aspects of the problem that you do not
master. On the contrary, if the approach you want to follow is
incremental by nature or known, EIC Pathfinder Open will not support
you.}
\end{quote}

\textbf{Question:} What logics are present?

\textbf{Analysis:} No primary coding on a conservative deductive pass.
The passage sets out an eligibility principle rooted in science logic
values --- genuine unknowns, intellectual risk, and the explicit
exclusion of incremental or known approaches. \enquote{Questions that
nobody knows the answer to yet} is a frontier-research criterion; the
rejection of \enquote{incremental by nature or known} is a science
logic value judgment. This is closer to a governing criterion than the
promotional register of [041]. However, the passage embodies science
logic values without enacting science logic governance --- it
encourages a scientific disposition but does not name who judges it or
on what formal evaluative basis. The evaluative machinery (peer
review, expert judgment, merit criteria) that would make
\texttt{LOGIC\_SCIENCE} unambiguous is absent. A close call; below
primary coding threshold on a conservative pass. Carried forward as an
observation for the inductive update.

\textbf{Coding:} None --- primarily promotional / eligibility framing.

\section*{Entry [044] --- EIC Pathfinder Open: Gatekeepers}

\textbf{Document location:} II.~EIC Pathfinder --- II.1~EIC Pathfinder
Open / Gatekeepers. \\
\textbf{Programme:} \texttt{[PATH]}.

\textbf{Source text:}
\begin{quote}
Before applying to this call, you should verify that your proposal
meets all the following essential characteristics
(\enquote{Gatekeepers}):
\begin{itemize}[leftmargin=*]
  \item Convincing long-term vision of a radically new technology that
    has the potential to have a transformative positive effect to
    solving a challenge in our economy and society.
  \item Concrete, novel and ambitious science-towards-technology
    breakthrough, providing advancement towards the envisioned
    technology.
  \item High-risk / high-gain research approach and methodology, with
    concrete and plausible objectives.
\end{itemize}
\end{quote}

\textbf{Question:} What logics are present?

\textbf{Analysis:} Opening (\enquote{Gatekeepers}) --- procedural. The
state is the implied background authority for any eligibility
criteria, but the introductory sentence does not argue a state
rationale. Does not rise to primary coding; consistent with the
principle established in [036]. First bullet ---
\texttt{LOGIC\_MARKET} + \texttt{LOGIC\_STATE}. \enquote{Economy}
invokes market logic; \enquote{society} and the possessive
\enquote{our} frame this as a European collective interest, giving it
a dual market / geopolitical dimension. The state is present in the
strategic framing of whose economy and whose society is being served.
Last two bullets --- \texttt{LOGIC\_SCIENCE}.
\enquote{Science-towards-technology breakthrough} and \enquote{research
approach and methodology} are explicit science logic criteria. No
\texttt{LOGIC\_TENSION} visible on the surface --- the three logics
are presented as complementary gatekeeping criteria rather than
pulling against each other.

\textbf{[?] Flag for inductive update:} Potential
\texttt{LOGIC\_TENSION} between the market / state outcome framing of
the first bullet and the science methodology framing of the last two
bullets. Whether the instrument is ultimately organised by what it
\emph{produces} (market disruption, societal challenge) or by how it
\emph{works} (scientific breakthrough, research methodology) may
reveal a tension not yet visible on a deductive pass.

\textbf{Coding:} \texttt{LOGIC\_SCIENCE} + \texttt{LOGIC\_MARKET} +
\texttt{LOGIC\_STATE}.

\section*{Entry [045] --- EIC Pathfinder Open: interdisciplinary research paragraph}

\textbf{Document location:} II.~EIC Pathfinder --- II.1~EIC Pathfinder
Open / interdisciplinary research. \\
\textbf{Programme:} \texttt{[PATH]}.

\textbf{Source text:}
\begin{quote}
\enquote{EIC Pathfinder Open involves interdisciplinary research and
development. By bringing diverse areas of research together, often
with different perspectives, terminologies and methodologies, within
individual projects and within a portfolio of projects, really new
things can be generated, and entirely new areas of research can be
opened up. It is up to you to compose the team that you need, that you
can learn from, and that you can move forward with.}
\end{quote}

\textbf{Question:} What logics are present?

\textbf{Analysis:} No primary coding on a conservative deductive pass.
The passage embodies science logic values throughout ---
interdisciplinary freedom, knowledge creation as the objective,
\enquote{entirely new areas of research can be opened up}, and
\enquote{it is up to you to compose the team} as an explicit statement
of researcher autonomy. Researcher autonomy is a core science logic
principle and is stated more explicitly here than in [043]. However,
consistent with [043], these are science logic values expressed in an
encouraging, applicant-facing register rather than science logic
governance. No peer authority, evaluative criteria, or merit framework
is invoked. The passage expresses permission rather than governing
criteria. Below primary coding threshold; strong \texttt{LOGIC\_SCIENCE}
undertone --- particularly researcher autonomy and knowledge creation
as objectives --- carried forward as an observation for the inductive
update.

\textbf{Coding:} None --- primarily promotional / applicant-facing.

\section*{Entry [046] --- EIC Pathfinder Open: expected outputs}

\textbf{Document location:} II.~EIC Pathfinder --- II.1~EIC Pathfinder
Open / expected outputs. \\
\textbf{Programme:} \texttt{[PATH]}.

\textbf{Source text:}
\begin{quote}
\enquote{The expected output of your project is the proof of principle
that the main ideas of the envisioned future technology are feasible,
thus validating its scientific and technological basis. Project
results should include top-level scientific publications in open
access. While your vision is expected to be worthwhile because of its
potential for future impact, for instance to create new markets,
improve our lives, or provide solutions for global challenges, these
are not expected to be achieved in the course of your EIC Pathfinder
Open project. However, you are expected to take the necessary measures
in the course of the project to allow future uptake to take place.
This includes: an adequate formal protection of the generated
Intellectual Property (IP), a plan for future exploitation and an
assessment of relevant aspects related to regulation, certification,
and standardisation.}
\end{quote}

\textbf{Question:} What logics are present?

\textbf{Analysis:} Both \texttt{LOGIC\_SCIENCE} and
\texttt{LOGIC\_MARKET} are present and clearly cross the primary
coding threshold. \texttt{LOGIC\_SCIENCE} --- strong and explicit.
\enquote{Proof of principle}, \enquote{validating its scientific and
technological basis}, and \enquote{top-level scientific publications
in open access} are core science logic deliverables. Science logic
governs what the project must produce now.
\texttt{LOGIC\_MARKET} --- explicit. \enquote{Adequate formal
protection of the generated IP}, \enquote{plan for future
exploitation}, and \enquote{assessment of relevant aspects related to
regulation, certification, and standardisation} are market-readiness
requirements embedded as mandatory deliverables within a science-phase
project. Market logic governs what must be prepared for later.
\texttt{LOGIC\_TENSION} --- visible on the surface of the text. The
passage creates an explicit temporal separation between the two
logics: \enquote{these are not expected to be achieved in the course
of your EIC Pathfinder Open project}. Science logic governs current
outputs; market logic governs mandatory future-facing preparations ---
both required simultaneously within the same project. The separation
is an institutional compromise between the two logics that is stated
rather than implied.

\textbf{Coding:} \texttt{LOGIC\_SCIENCE} + \texttt{LOGIC\_MARKET} +
\texttt{LOGIC\_TENSION}.

\section*{Note on the excerpt}

Entry [044] illustrates the conservative Pass~1 standard and the
\texttt{[?]} flagging convention; entries [043] and [045] illustrate
the threshold at which logic values are recognised but not coded;
entry [046] shows a fully coded passage where the tension between two
logics is stated in the text rather than inferred. Together the four
entries trace the discrimination on which the coding scheme depends.
The Pass~2 resolution of the [044] flag is recorded in
Appendix~\ref{app:codebook-inductive} (Pattern~A; Rule~1): the flagged
tension is constructed across bullets and so falls below the
surface-visibility threshold for \texttt{LOGIC\_TENSION}, but is
preserved as a Pattern~A observation.

